% !TEX root = chapter-standalone.tex

\section{Product of Categories}
\label{sec:product-of-categories}

\todotext{@J: review and edit}

Suppose we have two categories, each modeling a family of systems and processes, and we wish to consider the two \emph{side by side}, as a single composite system whose two parts do not interact.
The construction that achieves this is the \emph{product} of categories: an object of the product is a \emph{pair} of objects, one from each category; a morphism is a \emph{pair} of morphisms; and composition happens componentwise, each component minding its own business.

This is the categorical analogue of the product of \SY{monoids} (\cref{def:product-monoid}), and the definition follows the same componentwise pattern.

\begin{ctdefinition}[Cartesian product of categories]
    \label{def:cartesian-product-category}
    \SYNDEF{cartesian product of categories}
    Given two categories~$\CatC$ and~$\CatD$, their \emph{cartesian product}~$\CatC \Ctimes \CatD$ is the category specified as follows:
    \begin{enumerate}
        \item \emph{Objects}: Objects are pairs~$\tup{\Obja,\Objb}$, with~$\Obja\setin \ObC$ and~$\Objb\setin \ObD$.
        \item \emph{Morphisms}: A morphism from~$\tup{\Obja, \Objb}$ to~$\tup{\Objc,\Objd}$ is a pair of morphisms
              \begin{equation}\label{eq:prod-cat-morphisms}
                  \tup{\mora,\morb} \colon \tup{\Obja,\Objb}\mto \tup{\Objc,\Objd},
              \end{equation}
              with~$\mora\colon \Obja \mtoin\CatC \Objc$,~$\morb\colon \Objb\mtoin\CatD \Objd$.
        \item \emph{Composition}: The composition of morphisms is given by composing each component of the pair separately:
              \begin{equation}
                  \tup{\mora,\morb}\mthenof{\CatC\Ctimes\CatD} \tupp{\morc,\mord}=\tupp{\mora\mthenof{\CatC}\morc,\, \morb \mthenof{\CatD} \mord}.
              \end{equation}
        \item \emph{Identity morphisms}: Given objects~$\Obja\setin \ObC$ and~$\Objb\setin \ObD$, the \SY{identity morphism} on~$\tup{\Obja,\Objb}$ is the pair $\tup{\catidat\Obja,\catidat\Objb}$.
    \end{enumerate}
\end{ctdefinition}

\begin{remark}\label{rem:n-fold-cart-products-of-cats}
    In a manner analogous to \cref{def:cartesian-product-category} we can also define the cartesian product of any finite number of categories.
    For example, if $\CatA, \CatB, \CatC$ are three categories, their triple cartesian product $\CatA \Ctimes \CatB \Ctimes \CatC$ is a category whose objects are triples $\tup{\Obja, \Objb, \Objc}$, with $\Obja\setin \Ob_\CatA$, $\Objb\setin \Ob_\CatB$, $\Objc\setin \Ob_\CatC$.
    For any $n \setin \natnumbers$, the $n$-fold cartesian product of a category $\CatC$ with itself is denoted $\CatC^n$.
\end{remark}

\begin{exercise}
    \label{ex:product-cat-is-category}
    Check that \cref{def:cartesian-product-category} does define a category: verify that composition is associative and that the proposed \SY{identity morphisms} satisfy the identity laws.
\end{exercise}
\begin{solution}
    Both conditions can be checked componentwise, using the corresponding conditions in~\CatC and~\CatD.
    For associativity:
    \begin{equation}
        \begin{aligned}
            \pars{\tup{\mora, \morb} \mthen \tup{\morc, \mord}} \mthen \tup{\more, \morf}
            & = \tup{(\mora \mthenof\CatC \morc) \mthenof\CatC \more, \; (\morb \mthenof\CatD \mord) \mthenof\CatD \morf} \\
            & = \tup{\mora \mthenof\CatC (\morc \mthenof\CatC \more), \; \morb \mthenof\CatD (\mord \mthenof\CatD \morf)} \\
            & = \tup{\mora, \morb} \mthen \pars{\tup{\morc, \mord} \mthen \tup{\more, \morf}}.
        \end{aligned}
    \end{equation}
    For unitality: $\tup{\catidat\Obja, \catidat\Objb} \mthen \tup{\mora, \morb} = \tup{\catidat\Obja \mthenof\CatC \mora, \; \catidat\Objb \mthenof\CatD \morb} = \tup{\mora, \morb}$, and similarly on the other side.
\end{solution}

\subsection{Examples}

\begin{example}[Products of monoids, revisited]
    \label{exa:product-cat-monoids}
    For \SY{monoids}~$\monoidA$ and~$\monoidB$, viewed as one-object categories (\cref{def:monoidcat}), the product category is the one-object category of the product \SY{monoid}:
    \begin{equation}
        \monoidcat{\monoidA} \Ctimes \monoidcat{\monoidB}
        \; = \;
        \monoidcat{\pars{\monoidA \cartprod \monoidB}}.
    \end{equation}
    Indeed, an object of the left-hand side is a pair of the two single objects (so again a single object), and a morphism is a pair of monoid elements, composed componentwise---which is exactly composition in~$\monoidA \cartprod \monoidB$ (\cref{def:product-monoid}).
\end{example}

\begin{example}[Products of posets, revisited]
    \label{exa:product-cat-posets}
    Similarly, for \SY{posets}~$\posA$ and~$\posB$, viewed as categories (\cref{sec:posets-as-cats}), the product category~$\poscat{\posA} \Ctimes \poscat{\posB}$ is the category of the product poset (\cref{def:poset-product}): objects are pairs of elements, and there is a (single) morphism~$\tup{\posela, \poselb} \mto \tup{\posela', \poselb'}$ exactly when~$\posela \posAleq \posela'$ \emph{and}~$\poselb \posBleq \poselb'$---the componentwise order.
\end{example}

\begin{example}[Traveling with a wallet]
    \label{exa:product-cat-travel-money}
    Applied models compose well under products.
    Recall the category \Berg of hiking trips (\cref{sec:trekking}) and the currency category \Curr (\cref{sec:currency_cat}).
    An object of~$\Berg \Ctimes \Curr$ is a pair
    \begin{equation}
        \tup{\text{location}, \; \text{currency}},
    \end{equation}
    describing the state of a traveler: where they are, and in which currency they hold their money.
    A morphism is a pair (trip, exchange): a way of moving through the mountains \emph{together with} a way of converting one's money.
    Composition is componentwise: trips concatenate, exchanges compound.

    The product structure encodes the modeling assumption that the two aspects are \emph{independent}: which trail you hike does not constrain how you exchange money, and vice versa.
    If, in a finer model, the two aspects do interact (say, certain trips are only possible if one can pay a toll in the local currency), then the product category is no longer the right model---detecting such interaction is one reason to be explicit about when a product is being assumed.
\end{example}

\begin{remark}
    \label{rem:product-cat-projections-forward}
    Like the product of \SY{monoids} (\cref{rem:product-projections-forward}), the product of categories comes with two evident ``read off one component'' assignments, sending~$\tup{\Obja, \Objb}$ to~$\Obja$ (respectively~$\Objb$) and~$\tup{\mora, \morb}$ to~$\mora$ (respectively~$\morb$).
    Once we have introduced \SY{functors}, we will recognize these projections as functors, and see that they characterize the product of categories by a universal property.
\end{remark}

